\documentclass[11pt,a4paper]{article}
\input{../../sablony/pripravy/priprava.sty}
\usepackage{tikz}
\nastavPripravu{1. ročník}{Souhrnné opakování}{134}{Intervalová mapa: nerovnice a moduly}
\begin{document}
\blok{Úvod a připomenutí}
Opakování lineárních a kvadratických nerovnic, absolutní hodnoty, znaménkové tabulky, podílové nerovnice i průniku řešení. Studenti všechny výsledky znázorňují na číselné ose a rozlišují ostré a neostré nerovnosti.
\blok{Společně řešené příklady}
\begin{prikladReseniBox}[title={Příklad 1}]
Vyřešte $\dfrac{2x-3}{4}\le x+1$.
\tcblower
\[
\begin{aligned}
2x-3&\le4x+4\\
-7&\le2x\\
K&=\langle-\tfrac72;+\infty)
\end{aligned}
\]

\end{prikladReseniBox}
\begin{prikladReseniBox}[title={Příklad 2}]
Vyřešte $|2x-1|<3$.
\tcblower
\[
\begin{aligned}
-3&<2x-1<3\\
-2&<2x<4\\
K&=(-1;2)
\end{aligned}
\]

\end{prikladReseniBox}
\begin{prikladReseniBox}[title={Příklad 3}]
Vyřešte $x^2-5x+6\ge0$.
\tcblower
\[
\begin{aligned}
(x-2)(x-3)&\ge0\\
K&=(-\infty;2\rangle\cup\langle3;+\infty)
\end{aligned}
\]

\end{prikladReseniBox}
\begin{prikladReseniBox}[title={Příklad 4}]
Vyřešte $x^2-4|x|+3\le0$.
\tcblower
\[
\begin{aligned}
t&=|x|\ge0\\
(t-1)(t-3)&\le0\ \Rightarrow\ t\in\langle1;3\rangle\\
1\le |x|\le3&\ \Rightarrow\ K=\langle-3;-1\rangle\cup\langle1;3\rangle
\end{aligned}
\]

\end{prikladReseniBox}
\blok{Samostatné procvičování}
\begin{zadaniReseniBox}[title={Samostatná práce 1}]
Vyřešte $4-3x>1$.
\tcblower
\[
\begin{aligned}
-3x&>-3\\
K&=(-\infty;1)
\end{aligned}
\]

\end{zadaniReseniBox}
\begin{zadaniReseniBox}[title={Samostatná práce 2}]
Vyřešte $|x+2|\ge4$.
\tcblower
\[
\begin{aligned}
x+2&\le-4\ \text{nebo}\ x+2\ge4\\
K&=(-\infty;-6\rangle\cup\langle2;+\infty)
\end{aligned}
\]

\end{zadaniReseniBox}
\begin{zadaniReseniBox}[title={Samostatná práce 3}]
Vyřešte $\dfrac{(x-2)(x+1)}{x-3}>0$.
\tcblower
\[
\begin{aligned}
x&\ne3\\
\text{kritické body: }&-1,2,3\\
K&=(-1;2)\cup(3;+\infty)
\end{aligned}
\]

\end{zadaniReseniBox}
\begin{zadaniReseniBox}[title={Samostatná práce 4}]
Vyřešte zároveň $x^2-4\le0$ a $x>1$.
\tcblower
\[
\begin{aligned}
x^2-4&\le0\ \Rightarrow\ x\in\langle-2;2\rangle\\
K&=\langle-2;2\rangle\cap(1;+\infty)=(1;2\rangle
\end{aligned}
\]

\end{zadaniReseniBox}
\blok{Volitelná rezerva a výzva}
\poznamka{Výzva: vymysli nerovnici, jejíž řešení je právě $(-\infty;-2)\cup(3;+\infty)$. Jedna možnost: $(x+2)(x-3)>0$.}
\blok{Domácí úkol}
\begin{zadaniReseniBox}[title={Domácí úkol}]
\begin{enumerate}[label=\alph*),itemsep=2pt]
\item Vyřešte $-2x+3\le9$.
\item Vyřešte $|x-3|\le2$.
\item Vyřešte $x^2-x-6<0$.
\item Vyřešte $(x-4)(x+2)\ge0$.
\item Vyřešte $\dfrac{x-1}{x+2}\le0$.
\item Vyřešte $x^2-6|x|+8<0$.
\end{enumerate}
\tcblower
\begin{enumerate}[label=\alph*),itemsep=2pt]
\item
\[
\begin{aligned}
-2x&\le6\ \Rightarrow\ x\ge-3\\
K&=\langle-3;+\infty)
\end{aligned}
\]
\item
\[
\begin{aligned}
-2&\le x-3\le2\\
K&=\langle1;5\rangle
\end{aligned}
\]
\item
\[
\begin{aligned}
(x+2)(x-3)&<0\\
K&=(-2;3)
\end{aligned}
\]
\item
\[
\begin{aligned}
K&=(-\infty;-2\rangle\cup\langle4;+\infty)
\end{aligned}
\]
\item
\[
\begin{aligned}
x&\ne-2\\
\text{kritické body: }&-2,1\\
K&=(-2;1\rangle
\end{aligned}
\]
\item
\[
\begin{aligned}
t&=|x|\ge0\\
(t-2)(t-4)&<0\ \Rightarrow\ 2<t<4\\
K&=(-4;-2)\cup(2;4)
\end{aligned}
\]
\end{enumerate}
\end{zadaniReseniBox}
\poznamka{Čas: 5 + 19 + 16 + 5 minut. Nech studenty vedle algebraického zápisu skutečně nakreslit ose intervaly; podílová nerovnice je volitelná výzva pro rychlejší skupinu.}
\end{document}
